El avance en las tecnologías de automatización y robótica ha permitido el desarrollo de plataformas cada vez más accesibles y versátiles para la enseñanza y experimentación en entornos académicos. En este contexto, los robots manipuladores como el PhantomX Pincher, junto con herramientas de código abierto como ROS2 (Robot Operating System) y la Raspberry Pi 5, ofrecen un marco ideal para el diseño de sistemas inteligentes con aplicaciones prácticas en clasificación de objetos mediante visión por computador.

La visión artificial, combinada con técnicas de control robótico, permite a los sistemas identificar, categorizar y manipular objetos en su entorno. Esta capacidad resulta particularmente valiosa para procesos industriales, logísticos y de investigación. En el ámbito educativo, la integración de estas tecnologías proporciona a los estudiantes una plataforma para desarrollar competencias en programación, percepción computacional y control de sistemas embebidos.

El presente trabajo propone el diseño de un kit educativo que aproveche el potencial del robot PhantomX Pincher como plataforma base, incorporando una cámara, una Raspberry Pi 5, y el entorno ROS2 para permitir la clasificación de objetos a partir del procesamiento de imágenes en tiempo real. El enfoque del trabajo es conceptual y orientado al diseño, sin comprometerse con la construcción física del kit, y busca sentar las bases para futuras implementaciones prácticas o adaptaciones didácticas en cursos de robótica e inteligencia artificial.
